\documentclass[12pt,a4paper]{report}

\usepackage[utf8]{inputenc}
\usepackage[T1]{fontenc}
\usepackage[french]{babel}
\usepackage{lmodern}
\usepackage[top=2.5cm,bottom=2.5cm,left=3cm,right=2.5cm]{geometry}
\usepackage{setspace}
\onehalfspacing
\usepackage{parskip}
\setlength{\parindent}{0pt}

\usepackage{xcolor}
\usepackage{colortbl}
\definecolor{primary}{RGB}{30,80,110}
\definecolor{lightblue}{RGB}{220,235,245}
\definecolor{lightgray}{RGB}{245,245,245}
\definecolor{codebg}{RGB}{248,248,252}

\usepackage{booktabs}
\usepackage{tabularx}
\usepackage{array}
\newcolumntype{L}[1]{>{\raggedright\arraybackslash}p{#1}}
\newcolumntype{C}[1]{>{\centering\arraybackslash}p{#1}}

\usepackage{listings}
\lstset{
  backgroundcolor=\color{codebg},
  basicstyle=\footnotesize\ttfamily,
  keywordstyle=\color{primary}\bfseries,
  commentstyle=\color{gray}\itshape,
  stringstyle=\color{teal},
  numberstyle=\tiny\color{gray},
  numbers=left, numbersep=8pt,
  frame=single, rulecolor=\color{gray!30},
  breaklines=true, showstringspaces=false,
  language=Python,
  xleftmargin=15pt
}

\usepackage[hidelinks]{hyperref}
\usepackage{fancyhdr}
\pagestyle{fancy}
\fancyhf{}
\fancyhead[L]{\small\nouppercase{\leftmark}}
\fancyhead[R]{\small TER M1 --- Agentic AI}
\fancyfoot[C]{\thepage}
\renewcommand{\headrulewidth}{0.4pt}

\usepackage{enumitem}
\usepackage{amsmath}
\usepackage{float}
\usepackage{graphicx}

% Boîtes légères
\newcommand{\infobox}[2]{%
  \medskip
  \colorbox{lightblue}{\parbox{0.95\textwidth}{\smallskip\textbf{\color{primary}#1}\\[2pt]#2\smallskip}}
  \medskip
}

% ─────────────────────────────────────────────────────────────────────────────
\begin{document}

% ── Page de titre ─────────────────────────────────────────────────────────────
\begin{titlepage}
\begin{center}
\vspace*{1.5cm}
{\Large\bfseries Université de Montpellier --- Master Informatique}\\[0.4em]
{\large Intelligence Artificielle \& Data}\\[0.8em]
\rule{\textwidth}{1.5pt}\\[1.5em]
{\LARGE\bfseries\color{primary}
Architectures Logicielles pour l'IA Agentique\\[0.5em]}
{\large\bfseries
Systèmes Multi-Agents, Variabilité Architecturale\\
et Prototype de Recrutement Automatisé}\\[0.8em]
\rule{\textwidth}{1.5pt}\\[1.5em]
{\large Travaux d'Étude et de Recherche (TER) --- Master 1}\\[0.4em]
{\large Année universitaire 2025--2026}\\[2cm]
\begin{tabular}{ll}
\textbf{Auteurs :} & Sofiane Dzermane\\
& [Prénom Nom]\\[1em]
\textbf{Encadrant :} & M. Abdelhak-Djamel Seriai\\[1em]
\textbf{Date :} & Mai 2026\\
\end{tabular}
\end{center}
\end{titlepage}

% ── Résumé ────────────────────────────────────────────────────────────────────
\chapter*{Résumé}
\addcontentsline{toc}{chapter}{Résumé}

Ce rapport présente l'intégralité des travaux réalisés dans le cadre du TER de
Master~1 Informatique, portant sur l'étude des architectures logicielles pour
l'intelligence artificielle agentique. La démarche s'est articulée en trois temps :
une phase d'études bibliographiques approfondies sur les Systèmes Multi-Agents (SMA)
classiques et l'Agentic AI contemporaine, une phase de sélection du scénario applicatif,
et une phase de conception et développement d'un prototype opérationnel.

Le livrable final est un système de recrutement automatisé multi-agents implémentant
\textbf{huit agents spécialisés} orchestrés par LangGraph, une mémoire contextuelle
vectorielle inter-runs via ChromaDB (RAG), un mécanisme Human-in-the-Loop, deux modes
architecturaux déployables (monolithique et microservices Docker), une API REST avec
streaming SSE, une interface graphique NiceGUI et \textbf{91 tests automatisés}.

\medskip
\textbf{Mots-clés :} Systèmes Multi-Agents, Agentic AI, LangGraph, RAG, Groq,
Feature Model, Microservices, Human-in-the-Loop, SE4AI, AI4SE.

\chapter*{Abstract}
\addcontentsline{toc}{chapter}{Abstract}

This report presents the work carried out as part of the Master~1 Computer Science
research project (TER), focusing on the study and implementation of software
architectures for agentic AI-based systems. The methodology comprised a thorough
bibliographic study of classical Multi-Agent Systems and contemporary Agentic AI,
followed by the design and development of an automated recruitment system. The
prototype implements eight specialised LangGraph agents, cross-run vector memory
(ChromaDB), Human-in-the-Loop, Docker microservices, REST API with SSE streaming,
NiceGUI interface and 91 automated tests.

\tableofcontents
\listoffigures
\listoftables

% =============================================================================
\chapter{Introduction}
\label{chap:intro}
% =============================================================================

\section{Contexte général}

L'essor des modèles de langage de grande taille (LLM) depuis 2022 a profondément
reconfiguré la façon dont les systèmes logiciels intègrent l'intelligence artificielle.
Ces modèles ne se contentent plus de générer du texte : ils servent de \textit{couche
cognitive} dans des systèmes distribués capables de planifier, percevoir leur
environnement, sélectionner des outils et prendre des décisions autonomes. Ce mouvement,
désigné sous le terme d'\textbf{Agentic AI}, interroge fondamentalement les architectures
logicielles existantes.

Ce TER s'inscrit à l'intersection de deux axes complémentaires :

\begin{description}
  \item[AI4SE] (\textit{Artificial Intelligence for Software Engineering}) : l'IA
        comme outil au service de l'ingénierie logicielle --- génération de code,
        revue automatique, tests, optimisation architecturale.
  \item[SE4AI] (\textit{Software Engineering for AI}) : le génie logiciel comme
        cadre méthodologique pour industrialiser les systèmes IA --- architecture,
        gouvernance, traçabilité, tests, MLOps.
\end{description}

\section{Objectifs du TER}

L'encadrant, M. Seriai, a défini trois objectifs :

\begin{enumerate}
  \item \textbf{Étudier la variabilité architecturale} des systèmes logiciels basés
        sur l'IA et la formaliser sous forme d'un feature model.
  \item \textbf{Identifier des critères de choix} architecturaux orientant la
        sélection d'un style d'architecture adapté.
  \item \textbf{Implémenter une architecture de référence Agentic AI} : microservices,
        API Gateway, RAG, observabilité, Human-in-the-Loop, illustrée par un scénario
        applicatif réaliste.
\end{enumerate}

\section{Démarche globale}

\begin{enumerate}
  \item Phase bibliographique : 6 études SMA + 11 études Agentic AI + synthèses +
        analyse comparative.
  \item Phase de sélection du scénario : comparaison de plusieurs domaines d'application.
  \item Phase de développement itératif du prototype.
\end{enumerate}

% =============================================================================
\chapter{Phase d'études bibliographiques}
\label{chap:etudes}
% =============================================================================

\section{Inventaire du dossier d'études}

Le dossier \texttt{Etudes/} contient l'ensemble des productions de la phase
bibliographique. \textbf{17 articles scientifiques} ont été lus, analysés et
synthétisés. Trois documents de synthèse ont été produits.

\begin{table}[H]
\centering
\caption{Inventaire du dossier \texttt{Etudes/}}
\label{tab:etudes}
\begin{tabularx}{\textwidth}{L{5cm}C{1.5cm}L{6.5cm}}
\toprule
\textbf{Fichier} & \textbf{Type} & \textbf{Contenu} \\
\midrule
\texttt{etudes\_sma\_/etude\_1} & PDF & Cours fondateur SMA \\
\texttt{etudes\_sma\_/etude\_2} & PDF & Référence académique française \\
\texttt{etudes\_sma\_/etude\_3} & PDF & Visual survey agent-based computing \\
\texttt{etudes\_sma\_/etude\_4} & PDF & General survey théorie/applications \\
\texttt{etudes\_sma\_/etude\_5} & PDF & CMU : Multi-Agent Reinforcement Learning \\
\texttt{etudes\_sma\_/etude\_6} & PDF & Perspectives systèmes distribués \\
\texttt{etudes\_sma\_/synthese} & PDF & \textbf{Synthèse produite — SMA} \\
\addlinespace
\texttt{etudes\_agentic\_/etude\_1} & PDF & Wavestone : Agentic AI Playbook \\
\texttt{etudes\_agentic\_/etude\_2} & PDF & Abou Ali : double paradigme \\
\texttt{etudes\_agentic\_/etude\_3} & PDF & Alva \& Pandey : architecture modulaire \\
\texttt{etudes\_agentic\_/etude\_4} & PDF & Sapkota : AI Agents vs Agentic AI \\
\texttt{etudes\_agentic\_/etude\_5} & PDF & Patel : sécurité et risques LLM \\
\texttt{etudes\_agentic\_/etude\_6} & PDF & Hughes : safety et alignement \\
\texttt{etudes\_agentic\_/etude\_7} & PDF & Vaidhyanathan : comparaison frameworks \\
\texttt{etudes\_agentic\_/etude\_8} & PDF & TOSEM : SE4AI/AI4SE \\
\texttt{etudes\_agentic\_/etude\_9} & PDF & Milosevic : communautés agentiques \\
\texttt{etudes\_agentic\_/etude\_10} & PDF & Schneider : générative vers agentique \\
\texttt{etudes\_agentic\_/etude\_11} & PDF & Pati : paradigme Agentic AI \\
\texttt{etudes\_agentic\_/synthese} & PDF & \textbf{Synthèse produite --- Agentic AI} \\
\addlinespace
\texttt{analyse\_comparative\_SMA\_AgenticAI} & PDF & \textbf{Analyse croisée produite} \\
\texttt{comparaison\_frameworks.csv} & CSV & \textbf{Glossaire des concepts} \\
\bottomrule
\end{tabularx}
\end{table}

\section{Synthèse des études SMA}

\subsection{Définitions et fondements}

Selon Dorri et al. (2018), un Système Multi-Agents est \textit{``un ensemble d'agents
autonomes interagissant dans un environnement partagé pour atteindre des objectifs
individuels ou collectifs.''}
Maldonado et al. (2024) identifient cinq composants clés : Agents, Environnement,
Interactions, Organisation, Objectifs. Tout agent possède quatre propriétés
essentielles : \textbf{Autonomie}, \textbf{Réactivité}, \textbf{Proactivité},
\textbf{Sociabilité}.

\subsection{Architectures organisationnelles}

Trois styles sont distingués dans la littérature :

\begin{itemize}
  \item \textbf{Centralisée} : un agent maître coordonne les subordonnés. Avantage :
        cohérence globale. Inconvénient : point unique de défaillance.
  \item \textbf{Décentralisée} : la coordination émerge des interactions locales.
        Avantage : robustesse. Inconvénient : comportement global difficile à prédire.
  \item \textbf{Hiérarchique} : multi-niveaux, compromis le plus répandu en pratique.
\end{itemize}

\subsection{Protocoles formels de coordination}

Le \textit{General Survey} (étude 4) formalise : Contract Net Protocol, théorie des
jeux, mécanismes d'enchères, consensus distribué, et ACL (\textit{Agent Communication
Language}). Ces outils garantissent stabilité, optimalité et convergence.

\subsection{Le MARL : pont vers l'apprentissage}

Le CMU Paper (étude 5) établit un pont via le \textbf{Multi-Agent Reinforcement
Learning (MARL)}, permettant aux agents d'adapter dynamiquement leurs stratégies.
Les défis : non-stationnarité, instabilité de convergence, équilibre
coopération/compétition.

\section{Synthèse des études Agentic AI}

\subsection{Clarification conceptuelle}

La synthèse Agentic AI converge vers une clarification essentielle : l'Agentic AI
ne se réduit ni à un agent LLM augmenté d'outils, ni aux SMA classiques. Trois
niveaux distincts émergent :

\begin{description}
  \item[AI Agents] (Sapkota et al.) : systèmes modulaires LLM, spécialisés via
        prompt engineering.
  \item[Agentic AI] (Schneider ; Pati) : paradigme orienté objectifs, planification
        autonome, mémoire persistante, adaptation dynamique.
  \item[Communautés agentiques] (Milosevic et Rabhi) : écosystèmes hybrides
        humains-IA avec gouvernance formelle.
\end{description}

La trajectoire évolutive identifiée est :
\[
\text{GenAI} \to \text{AI Agents} \to \text{Agentic AI} \to
\text{Communautés agentiques} \to \text{Écosystèmes hybrides gouvernés}
\]

\subsection{Double paradigme : symbolique vs neuronal}

Abou Ali et al. (2026) proposent un cadre analytique fondamental :

\begin{table}[H]
\centering
\caption{Double paradigme Symbolique vs Neuronal (Abou Ali et al., 2026)}
\label{tab:paradigmes}
\begin{tabularx}{\textwidth}{L{3cm}L{5.5cm}L{5cm}}
\toprule
\textbf{Paradigme} & \textbf{Caractéristiques} & \textbf{Domaines} \\
\midrule
\textbf{Symbolique} & Planification déterministe, état persistant, vérification formelle & Santé critique \\
\textbf{Neuronal} & Orchestration LLM, génération probabiliste, adaptabilité élevée & Finance, environnements dynamiques \\
\bottomrule
\end{tabularx}
\end{table}

Les études convergent vers une trajectoire \textbf{neuro-symbolique hybride}, que
notre prototype incarne directement : agents LLM (neuronal) + guards déterministes
(symbolique).

\subsection{Comparaison des frameworks}

L'étude 7 (Vaidhyanathan et Taibi, 2026) et le glossaire CSV produit constituent
la pièce maîtresse de la sélection technique :

\begin{table}[H]
\centering
\caption{Comparaison des frameworks Agentic AI}
\label{tab:frameworks}
\begin{tabularx}{\textwidth}{L{2.3cm}L{2.5cm}C{1.8cm}C{1.5cm}L{3.8cm}}
\toprule
\textbf{Framework} & \textbf{Coordination} & \textbf{État partagé} & \textbf{HITL} & \textbf{Points forts} \\
\midrule
\rowcolor{lightblue}
\textbf{LangGraph} & Graphe orienté & TypedDict riche & Natif & Contrôle fin, traçabilité \\
CrewAI & Rôles/délégation & Partiel & Limité & Prise en main rapide \\
AutoGen & Conversation & Limité & Oui & Flexibilité, multi-LLM \\
OpenAI Swarm & Handoff & Absent & Non & Légèreté \\
MetaGPT & SOP rigides & Oui & Partiel & Processus logiciel \\
\bottomrule
\end{tabularx}
\end{table}

LangGraph est le seul framework qui supporte nativement le blackboard (TypedDict),
le fan-out parallèle (\texttt{Send()}), le HITL sans contournement, et est
indépendant du fournisseur LLM.

\section{Analyse comparative SMA / Agentic AI}

Le document \texttt{analyse\_comparative\_SMA\_AgenticAI.pdf} identifie cinq
convergences et cinq divergences structurelles.

\medskip
\colorbox{lightblue}{\parbox{0.95\textwidth}{\smallskip
\textbf{\color{primary}Conclusion majeure de l'analyse comparative}\\
\textit{``Les SMA ne s'opposent pas à l'Agentic AI : ils constituent un réservoir
de solutions formelles et éprouvées capable de combler ses faiblesses actuelles
en matière de stabilité, de gouvernance et d'explicabilité.''}
\smallskip}}
\medskip

\begin{table}[H]
\centering
\caption{Divergences structurelles SMA vs Agentic AI}
\label{tab:divergences}
\begin{tabularx}{\textwidth}{L{3cm}L{5.5cm}L{4.5cm}}
\toprule
\textbf{Critère} & \textbf{SMA classiques} & \textbf{Agentic AI} \\
\midrule
Raisonnement & Logique, déterministe, prouvé & Probabiliste (LLM) \\
Coordination & Protocoles formels (ACL, Contract Net) & Prompt engineering \\
Apprentissage & Optionnel (MARL) & Central et indispensable \\
Maturité & Décennies, bases mathématiques & Écosystème fragmenté \\
Explicabilité & Règles traçables & Déficit inhérent aux LLM \\
\bottomrule
\end{tabularx}
\end{table}

% =============================================================================
\chapter{Exploration des sujets et choix du scénario}
\label{chap:sujets}
% =============================================================================

\section{Sujets envisagés}

\subsection{Scraping et comparateur de prix}

\textbf{Concept :} Agents autonomes parcourant des sites e-commerce, collectant
des prix et alertant sur les meilleures offres.

\textbf{Architecture SMA envisagée :} Agent Scraper ($\times N$, un par site),
Agent Normalisateur, Agent Comparateur, Agent Alerteur.

\textbf{Pourquoi non retenu :} La valeur ajoutée du LLM est marginale --- la
comparaison de prix est un problème algorithmique pur. Le pipeline est trop linéaire
pour justifier des patterns SMA complexes. Il manque la dimension cognitive de
l'Agentic AI.

\subsection{Assistant de veille technologique}

\textbf{Concept :} Agents spécialisés par domaine (IA, cybersécurité, cloud) qui
agrègent, filtrent et synthétisent des informations depuis RSS, arXiv, GitHub Trending.

\textbf{Pourquoi non retenu :} Périmètre difficile à délimiter. L'évaluation
qualitative des résultats est subjective. L'encadrant a orienté vers un sujet avec
des métriques plus objectives.

\subsection{Système de recommandation personnalisée}

\textbf{Concept :} Plusieurs agents modélisant différents aspects du profil
utilisateur (historique, préférences, comportement), convergeant par vote pondéré.

\textbf{Pourquoi non retenu :} Nécessite des données utilisateur historiques
difficiles à simuler de façon réaliste. Les métriques de qualité demandent un jeu
de données conséquent.

\section{Le rôle de l'encadrant dans le choix final}

Lors d'une séance de suivi, \textbf{M. Seriai} a proposé le sujet du recrutement
automatisé. Cette suggestion a immédiatement convergé les réflexions :

\begin{enumerate}
  \item Le recrutement est un \textbf{pipeline naturellement multi-étapes} : analyser,
        chercher, filtrer, évaluer, vérifier, contacter. Chaque étape = un agent avec
        une responsabilité unique (principe SRP appliqué aux agents).
  \item Le \textbf{HITL y est parfaitement naturel} : décider de contacter un candidat
        est une décision à forte valeur humaine.
  \item Le \textbf{RAG est démontrable} : les profils invalides récurrents entre runs
        justifient une blacklist inter-runs avec valeur mesurable.
  \item Les \textbf{métriques sont objectives} : scores de matching, durée de run.
\end{enumerate}

\section{Critères de décision comparés}

\begin{table}[H]
\centering
\caption{Comparaison des sujets envisagés}
\label{tab:sujets}
\begin{tabularx}{\textwidth}{L{3.5cm}C{2cm}C{2cm}C{2cm}C{2cm}}
\toprule
\textbf{Critère} & \textbf{Prix} & \textbf{Veille} & \textbf{Reco.} & \textbf{Recrutement} \\
\midrule
Réalisme professionnel & Moyen & Élevé & Moyen & \textbf{Très élevé} \\
Complexité SMA justifiée & Faible & Moyen & Moyen & \textbf{Élevée} \\
HITL naturel & Non & Non & Non & \textbf{Oui} \\
RAG démontrable & Non & Oui & Oui & \textbf{Oui} \\
Métriques objectives & Oui & Non & Oui & \textbf{Oui} \\
Faisabilité gratuite & Oui & Oui & Non & \textbf{Oui} \\
\bottomrule
\end{tabularx}
\end{table}

% =============================================================================
\chapter{Conception et architecture du prototype}
\label{chap:conception}
% =============================================================================

\section{Vision générale}

Le système automatise le sourcing de candidats à partir d'une fiche de poste en
langage naturel. L'entrée est une chaîne de caractères (ex : \textit{``Développeur
Python senior, 5 ans, Paris, FastAPI, LangGraph, Docker, RAG''}). La sortie est un
rapport structuré avec candidats classés par score (0--100), statut de validation,
liens de profils et messages de contact personnalisés.

\section{Schéma logique du pipeline}

\begin{figure}[H]
\centering
\begin{tabularx}{\textwidth}{X}
\toprule
\textbf{Pipeline multi-agents LangGraph} \\
\midrule
\texttt{START} $\to$ \textbf{A1 Orchestrateur} $\to$ \textbf{A2 Analyste}
$\to$ \textbf{A3a Stratège} \\
$\to$ \textbf{A3b Collecteur} \textit{(DuckDuckGo + GitHub API + Stack Overflow)} \\
$\to$ \textbf{A3c Filtre} \textit{(30+ domaines agrégateurs, 50+ keywords, blacklist RAG)} \\
$\to$ \textbf{A6 Déduplicateur} \\
$\xrightarrow{\text{Send()}\times N}$ \textbf{A4 Évaluateur} \textit{(Groq LLM, cache RAG)} $\to$ \texttt{reduce\_scores} \\
$\xrightarrow{\text{Send()}\times N}$ \textbf{A5 Vérificateur} \textit{(LLM + guards déterministes)} $\to$ \texttt{reduce\_validations} \\
$\to$ \textbf{Injection RAG} \textit{(ChromaDB, cosine $\geq 0.75$, récents $\leq 30$ jours)} \\
$\xrightarrow{\text{conditionnel}}$ \textbf{A7 Recruteur} $\to$ \textbf{Rapport} $\to$ \textbf{A8 Persistance} $\to$ \texttt{END} \\
\bottomrule
\end{tabularx}
\caption{Architecture logique du pipeline multi-agents}
\label{fig:pipeline}
\end{figure}

\section{Le blackboard : cœur formel du SMA}

Le fichier \texttt{src/state.py} définit le blackboard, pierre angulaire de
l'architecture SMA. C'est un \texttt{TypedDict} partagé entre tous les agents, avec
une \textbf{convention d'écriture stricte} : chaque agent possède ses champs dédiés.

\begin{lstlisting}[caption={src/state.py --- Blackboard TypedDict avec reducers}]
class GraphState(TypedDict):
    fiche_poste: str                    # Input entreprise
    profil_competences: dict            # A2 ecrit, lu par A3a, A4
    requetes_recherche: dict            # A3a ecrit, lu par A3b
    resultats_bruts: list[dict]         # A3b ecrit, lu par A3c
    # Annotated + operator.add = reducer pour fan-out parallele
    profils_bruts: Annotated[list[Candidat], operator.add]
    profils_dedupliques: list[Candidat] # A6 ecrit, lu par A4(xN)
    candidats_scores: Annotated[list[CandidatScore], operator.add]
    candidats_valides: Annotated[list[CandidatValide], operator.add]
    messages_envoyes: Annotated[list[dict], operator.add]
    rapport_final: str                  # A1 ecrit en dernier
\end{lstlisting}

Les champs \texttt{Annotated[list, operator.add]} sont des \textbf{reducers} :
lors du fan-out \texttt{Send()}, N instances écrivent simultanément dans ces listes
et LangGraph les fusionne automatiquement, sans verrou explicite. C'est le pattern
\textbf{Map-Reduce} appliqué à un SMA.

\section{Construction du graphe}

\begin{lstlisting}[caption={src/graph.py --- Fan-out et HITL}]
def route_to_evaluateurs(state: GraphState) -> list[Send]:
    """Fan-out : 1 instance A4 par profil, execution parallele."""
    return [
        Send("evaluateur", {
            "candidat": candidat,
            "profil_competences": state["profil_competences"],
        })
        for candidat in profils[:MAX_PROFILS_PARALLELES]
    ]

# Compilation avec HITL
app = graph.compile(
    checkpointer=MemorySaver(),
    interrupt_before=["recruteur"]   # pause avant A7
)
\end{lstlisting}

\section{Feature model --- Variabilité architecturale}

\begin{figure}[H]
\centering
\begin{tabularx}{\textwidth}{L{3.2cm}C{1.5cm}L{8cm}}
\toprule
\textbf{Feature} & \textbf{Type} & \textbf{Variantes} \\
\midrule
\rowcolor{lightblue}
\textbf{SMA Recrutement} & racine & --- \\
Mode d'exécution & XOR & \textbf{Local} (monolithique, LangGraph) $|$ \textbf{Microservices} (Docker) \\
Fournisseur LLM & mand. & \textbf{Groq} (\texttt{llama-3.3-70b-versatile}) [extensible] \\
HITL & XOR & \textbf{Activé} \texttt{interrupt\_before=["recruteur"]} $|$ \textbf{Désactivé} \\
Interface & OR & \textbf{CLI} $|$ \textbf{API REST} $|$ \textbf{UI Web} \\
Mémoire RAG & AND & \textbf{Blacklist} inter-runs $+$ \textbf{Cache} de scores \\
\bottomrule
\multicolumn{3}{l}{\scriptsize mand.~= mandatoire \quad XOR~= alternatif exclusif \quad OR~= un ou plusieurs \quad AND~= tous obligatoires}
\end{tabularx}
\caption{Feature model du système de recrutement multi-agents}
\label{fig:featuremodel}
\end{figure}

\section{Patterns SMA implémentés}

\begin{table}[H]
\centering
\caption{Patterns SMA implémentés dans le prototype}
\label{tab:patterns}
\begin{tabularx}{\textwidth}{L{3.5cm}L{4cm}L{6cm}}
\toprule
\textbf{Pattern SMA} & \textbf{Implémentation} & \textbf{Justification} \\
\midrule
Superviseur & A1 via flux du graphe & Topologie étoile, délégation formelle \\
Blackboard & \texttt{GraphState} TypedDict & État partagé, droits séparés par agent \\
Send / Map-Reduce & A4+A5 via \texttt{Send()} & $N$ instances parallèles \\
Validation pair-à-pair & A4 score $\to$ A5 contrôle & Séparation évaluation/vérification \\
Routage conditionnel & \texttt{route\_apres\_verification()} & Décision sur meilleur score \\
Human-in-the-Loop & \texttt{interrupt\_before=["recruteur"]} & Suspension pour validation \\
\bottomrule
\end{tabularx}
\end{table}

\section{Justification des choix technologiques}

\begin{table}[H]
\centering
\caption{Justification des choix technologiques}
\label{tab:technos}
\begin{tabularx}{\textwidth}{L{2.5cm}L{2.5cm}L{8cm}}
\toprule
\textbf{Techno} & \textbf{Alternative} & \textbf{Justification} \\
\midrule
LangGraph & CrewAI, AutoGen & Graphe = SMA formel, état riche, fan-out, HITL natif, pas de vendor lock-in \\
Groq (cloud) & Ollama (local) & Ollama crashait la machine (8~Go RAM), Groq free tier 100k tok/jour \\
ChromaDB & Pinecone & Local, gratuit, sans serveur externe \\
FastAPI & Flask & Async natif, SSE simplifié, swagger auto-généré \\
DuckDuckGo & SerpAPI & Gratuit, sans clé API \\
LangGraph & Kafka & LangGraph \textit{est} le bus : arête = canal, state = message \\
\bottomrule
\end{tabularx}
\end{table}

% =============================================================================
\chapter{Implémentation détaillée}
\label{chap:impl}
% =============================================================================

\section{Structure du projet}

\begin{description}[font=\ttfamily]
  \item[src/agents/] 8 agents + nœud injection RAG ($\sim$1\,400 lignes)
  \item[src/tools/] RAG, scraping, recherche DDG, GitHub API, Stack Overflow API
  \item[src/ui/] Interface NiceGUI (3 fichiers : app, layout, api\_client)
  \item[src/api.py] API Gateway FastAPI + SSE + SQLite ($\sim$500 lignes)
  \item[src/graph.py] Construction du graphe LangGraph
  \item[src/state.py] Blackboard TypedDict avec reducers
  \item[src/config.py] Configuration + \texttt{get\_llm()} avec retry
  \item[src/prompts.py] 4 prompts système centralisés
  \item[src/observabilite.py] Métriques par nœud, export JSON
  \item[services/] 6 microservices HTTP FastAPI (mode Docker)
  \item[tests/] 91 tests automatisés
\end{description}

\section{Les huit agents}

\subsection{A1 --- Orchestrateur (Superviseur)}

Point d'entrée et sortie du pipeline. Initialise le run, produit le rapport final
Markdown. Implémente le pattern Superviseur : délégation formelle via le flux du
graphe.

\subsection{A2 --- Analyste de poste}

Transforme la fiche de poste en JSON structuré via Groq. Prompt system avec
règle explicite : \textit{``AWS, GCP, Azure ne sont PAS des villes.''} (bug
identifié et corrigé en production). Enrichissement déterministe : fonction
\texttt{\_inferer\_niveau\_experience()} avec regex pour les patterns alternance,
stage, junior, senior.

\textbf{Exemple d'entrée/sortie :}

\medskip
\begin{tabularx}{\textwidth}{L{5.5cm}L{8cm}}
\textit{Entrée :} \texttt{"Développeur Python senior, 5 ans, Paris, FastAPI, LangGraph"} &
\textit{Sortie JSON :} \texttt{\{"hard\_skills": ["Python","FastAPI","LangGraph"], "niveau\_experience": "senior", "experience\_min": 5, "localisations": ["Paris"]\}} \\
\end{tabularx}

\subsection{A3 --- Division en trois sous-agents}

L'agent collecteur initial a été \textbf{refactorisé en A3a/A3b/A3c} lors d'une
itération architecturale, car la responsabilité unique était violée.

\paragraph{A3a --- Stratège.} Génère un plan de recherche par source via LLM.
Opérateurs déterministes : \texttt{site:linkedin.com/in/}, \texttt{language:python
followers:>5}.

\paragraph{A3b --- Collecteur.} Exécute le plan : DuckDuckGo, GitHub API officielle,
Stack Overflow API. Délai 1,5~s entre appels DDG (rate limit). Déduplication par URL.

\paragraph{A3c --- Filtre anti-bruit.} Filtre algorithmique pur en 4 passes :

\begin{lstlisting}[caption={src/agents/chercheur\_filtre.py --- Blacklist RAG}]
# Passe 3 : blacklist RAG (avant scraping = avant depenser des appels)
if _memoire is not None and _memoire.est_blackliste(url):
    n_blacklist_drop += 1
    continue  # URL invalide connue, skip immediat

# Passe 4 : post-filtre apres scraping
if _is_noise_text(profil_brut) or _is_aggregated_profile(title, profil_brut):
    n_post_drop += 1
    continue
\end{lstlisting}

\textbf{Résultats sur le run de référence :} 46 hits bruts $\to$ 13 après filtrage
(taux de bruit : 68\,\%).

\subsection{A6 --- Déduplicateur}

Fusionne les profils identiques par URL ou similarité de nom ($>85\,\%$ via
\texttt{SequenceMatcher}). En cas de fusion, les \texttt{profil\_brut} sont
concaténés pour maximiser l'information disponible à A4.

\subsection{A4 --- Évaluateur ($\times N$ parallèles, avec cache RAG)}

Évalue UN candidat par rapport au profil de compétences. Barème strict dans le
prompt : 85--100 (excellent), 70--84 (bon), 50--69 (intéressant), 30--49 (partiel),
0--29 (hors sujet). Cache RAG avant appel LLM :

\begin{lstlisting}[caption={src/agents/evaluateur.py --- Cache RAG}]
cache = get_memoire().get_score_cache(url, fiche_poste)
if cache is not None:
    # Score reutilise sans appel LLM => economie 1 appel Groq
    return {"candidats_scores": [CandidatScore(score_global=cache["score"])]}
\end{lstlisting}

\subsection{A5 --- Vérificateur ($\times N$ parallèles) --- La pièce de résistance}

Validation pair-à-pair en deux couches combinées :

\textbf{Couche LLM :} contrôle la cohérence sémantique du score A4, détecte les
incohérences, valide/invalide/doute.

\textbf{Guards déterministes :}

\begin{lstlisting}[caption={src/agents/verificateur.py --- Guards anti-faux positifs}]
# Guard 1 : page non-candidat (offre emploi, formation, article...)
if _profil_non_candidat(nom, profil_brut, url):
    score_final = min(score_final, 20.0)
    statut = "invalide"

# Guard 2 : adequation minimale (hard skills visibles dans le texte)
ok, raison = _adequation_minimale(profil_requis, nom, profil_brut)
if not ok:
    score_final = min(score_final, 45.0); statut = "invalide"

# Guard 3 : experience incompatible (senior pour un poste junior, etc.)
incompatible, raison = _experience_incompatible(profil_brut, niveau_requis)
if incompatible:
    score_final = min(score_final, 45.0); statut = "invalide"

# Guard 4 : upgrade LinkedIn (scraping bloque, si score A4 >= 75 on fait confiance)
if statut == "douteux" and score_a4 >= SCORE_SEUIL_CONTACT and "/in/" in url:
    statut = "valide"  # LinkedIn bloque le scraping : A5 voit trop peu de contenu
\end{lstlisting}

Sur le run de référence, \textbf{3 candidats valides sur 5} ont bénéficié de
l'upgrade automatique LinkedIn.

\subsection{Nœud Injection RAG}

Après la réduction des validations A5, enrichit les résultats avec des candidats
connus des runs précédents pour une fiche similaire (similarité cosine $\geq 0.75$,
$\leq 30$ jours). Anti-doublon par URL.

\subsection{A7 --- Recruteur}

Rédige des messages de contact personnalisés (max 150 mots). Mode absolu : score
$\geq 75$. Mode relatif : top-3 parmi les viables si aucun n'atteint le seuil.

\subsection{A8 --- Persistance RAG}

Nœud terminal. Stratégie de persistance :
\begin{itemize}[noitemsep]
  \item \texttt{"valide"} $\to$ cache de score pour runs futurs (A4)
  \item \texttt{"invalide"} $\to$ blacklist inter-runs (A3c)
  \item \texttt{"douteux"} $\to$ \textbf{non persisté} (incertitude trop élevée)
\end{itemize}

Règle : le statut ``invalide'' (rejet HITL) l'emporte sur ``valide'' lors de la
déduplication par \texttt{candidat\_id}.

\section{Architecture de la mémoire RAG}

ChromaDB local (\texttt{./data/chromadb/}) + SentenceTransformers
(\texttt{all-MiniLM-L6-v2}, $\sim$80~Mo). Deux collections :

\begin{itemize}
  \item \texttt{fiches\_poste} : embeddings des fiches, id = hash SHA1.
  \item \texttt{candidats\_evalues} : profils + métadonnées (score, statut, url,
        fiche\_id, last\_seen).
\end{itemize}

Le \texttt{fiche\_id} rattache chaque candidat à la fiche pour laquelle il a été
évalué : sans ce lien, un candidat Python serait biaisé pour un poste Java.

\medskip
\colorbox{lightblue}{\parbox{0.95\textwidth}{\smallskip
\textbf{\color{primary}Propriété d'apprentissage cumulatif}\\
Le système accumule de la connaissance opérationnelle sans modifier les poids du
modèle. Chaque run enrichit la base et rend les suivants plus rapides et plus précis.
Sur le Run~2 d'une même fiche : gain estimé de 30--40\,\% de durée (blacklist +
cache + injection).
\smallskip}}

\section{Human-in-the-Loop}

L'API expose trois actions via \texttt{POST /runs/\{id\}/hitl} :
\begin{description}[noitemsep]
  \item[\texttt{approve}] Reprend le pipeline tel quel.
  \item[\texttt{skip}] Invalide tous les candidats (mémorisé en RAG via A8).
  \item[\texttt{edit}] Remplace la liste de candidats par une version JSON éditée.
\end{description}

\section{Mode microservices}

\begin{table}[H]
\centering
\caption{Services Docker du mode microservices}
\label{tab:microservices}
\begin{tabularx}{\textwidth}{L{3.5cm}C{1.8cm}L{7.5cm}}
\toprule
\textbf{Service} & \textbf{Port} & \textbf{Rôle} \\
\midrule
\texttt{svc-analyste} & 8001 & A2 : analyse fiche de poste \\
\texttt{svc-chercheur} & 8002 & A3a/b/c : collecte + blacklist RAG \\
\texttt{svc-evaluateur} & 8003 & A4 : scoring + cache RAG \\
\texttt{svc-verificateur} & 8004 & A5 : vérification + guards déterministes \\
\texttt{svc-recruteur} & 8005 & A7 : rédaction messages \\
\texttt{svc-orchestrateur} & 8006 & Orchestration HTTP (remplace LangGraph) \\
\texttt{api-gateway} & 8000 & Point d'entrée public FastAPI \\
\texttt{sma-ui} & 8080 & Interface NiceGUI \\
\bottomrule
\end{tabularx}
\end{table}

\begin{table}[H]
\centering
\caption{Trade-offs monolithique vs microservices}
\label{tab:tradeoffs}
\begin{tabularx}{\textwidth}{L{3cm}L{5.5cm}L{5cm}}
\toprule
\textbf{Critère} & \textbf{Monolithique (LangGraph)} & \textbf{Microservices (Docker)} \\
\midrule
Latence & Optimale (in-process) & +30--50\,\% (HTTP) \\
Déployabilité & 1 processus & Cloud-ready \\
Scaling & Vertical & Horizontal par service \\
Débogage & Traces LangGraph & Logs distribués \\
RAG & Blacklist + cache + injection & Blacklist + cache \\
\bottomrule
\end{tabularx}
\end{table}

\section{Tests automatisés}

\begin{table}[H]
\centering
\caption{Couverture des 91 tests}
\label{tab:tests}
\begin{tabularx}{\textwidth}{L{6cm}C{1.8cm}L{6cm}}
\toprule
\textbf{Fichier} & \textbf{Tests} & \textbf{Couverture} \\
\midrule
\texttt{test\_graph.py} & 28 & Topologie, 16 nœuds, arêtes \\
\texttt{test\_verificateur\_rules.py} & 5 & Guards déterministes A5 \\
\texttt{test\_stratege\_queries.py} & 15 & Requêtes A3a \\
\texttt{test\_api.py} & 43 & Endpoints, HITL, SQLite \\
\midrule
\textbf{Total} & \textbf{91} & \textbf{91/91 passent} \\
\bottomrule
\end{tabularx}
\end{table}

% =============================================================================
\chapter{Évaluation et résultats}
\label{chap:eval}
% =============================================================================

\section{Run de référence}

Fiche : \textit{``Développeur Python senior, 5~ans, Paris ou remote, FastAPI,
LangGraph, Docker, RAG''} (21 mai 2026, \texttt{llama-3.3-70b-versatile}).

\begin{table}[H]
\centering
\caption{Métriques du run de référence}
\label{tab:metriques}
\begin{tabularx}{\textwidth}{L{5.5cm}C{2.5cm}L{5cm}}
\toprule
\textbf{Étape} & \textbf{Valeur} & \textbf{Commentaire} \\
\midrule
Profils bruts collectés (A3b) & 46 & DDG + GitHub + SO \\
Après filtrage A3c & 13 & Taux de bruit : 68\,\% \\
Après déduplication A6 & 12 & 1 doublon nom \\
Candidats évalués A4 & 10 & Limite MAX\_PROFILS=10 \\
Candidats valides A5 & 5 & Dont 3 upgrades LinkedIn \\
Messages rédigés A7 & 2 & Mode absolu (score $\geq 75$) \\
\textbf{Durée totale} & \textbf{173~s} & 10 appels LLM parallèles \\
\bottomrule
\end{tabularx}
\end{table}

\begin{table}[H]
\centering
\caption{Meilleurs candidats identifiés}
\label{tab:candidats}
\begin{tabularx}{\textwidth}{L{4cm}C{1.5cm}C{2cm}L{5.5cm}}
\toprule
\textbf{Candidat} & \textbf{Score} & \textbf{Source} & \textbf{Remarques A5} \\
\midrule
Srikanth Vanjre & 88 & LinkedIn & Senior Backend AI/RAG --- excellent \\
Shubhangi Gaur & 80 & LinkedIn & Python/FastAPI senior \\
Nikita Koznev & 72 & LinkedIn & Python senior (upgrade auto) \\
Jyotsna Konchada & 68 & LinkedIn & Backend Python (upgrade auto) \\
Vamshi Krishna & 68 & LinkedIn & Full Stack Python (upgrade auto) \\
\bottomrule
\end{tabularx}
\end{table}

\section{Apport du RAG sur les runs répétés}

\medskip
\colorbox{lightblue}{\parbox{0.95\textwidth}{\smallskip
\textbf{\color{primary}Comparaison Run~1 vs Run~2 --- même fiche de poste}\\[4pt]
\textbf{Run~1 (base RAG vide) :} 0 URL blacklistée, 0 injection, durée 173~s.\\[3pt]
\textbf{Run~2 (base RAG peuplée) :}
Les profils SO invalides (Jon Skeet C\#, VonC...) sont filtrés \textit{avant
scraping} par A3c : économie de 3 scrapings + 3 appels A4 + 3 appels A5.
Les 5 candidats valides du Run~1 sont réinjectés automatiquement.
Gain estimé : $-$6 appels Groq, durée réduite de 30--40\,\%.
\smallskip}}
\medskip

\section{Analyse critique}

\paragraph{Points forts.}
Architecture formellement fondée (6 patterns SMA), guards déterministes, RAG cumulatif
inter-runs, observabilité structurée, dual mode (monolithique + microservices), 91 tests.

\paragraph{Limites assumées.}
Scraping LinkedIn bloqué à 100\,\% (upgrade A5 comme palliatif) ; quota Groq
100k~tokens/jour ($\sim$3--5 runs complets) ; calibration empirique des seuils
RAG (0,75/0,85) sur un faible nombre de runs ; injection RAG absente en mode
microservices (blacklist et cache seuls actifs).

% =============================================================================
\chapter{Conclusion et perspectives}
\label{chap:conclusion}
% =============================================================================

\section{Bilan}

Ce TER a atteint les trois objectifs assignés :

\begin{enumerate}
  \item \textbf{État de l'art structuré} : 17 articles analysés, 4 documents de
        synthèse produits, comparaison de 5 frameworks, trajectoire évolutive
        identifiée (GenAI $\to$ Communautés agentiques).
  \item \textbf{Feature model} : variabilité en 5 dimensions formalisée, décisions
        architecturales justifiées par des contraintes métier mesurables.
  \item \textbf{Prototype opérationnel} : 8 agents LangGraph, RAG inter-runs
        (blacklist + cache + injection), HITL, microservices Docker, API SSE,
        UI NiceGUI, 91 tests --- $\sim$173~s end-to-end.
\end{enumerate}

Ce prototype démontre empiriquement que l'Agentic AI, structurée par les patterns
formels des SMA (blackboard, superviseur, map-reduce, validation pair-à-pair),
produit un système fiable, traçable et évolutif --- ce que les implémentations naïves
de LLM enchaînés ne garantissent pas.

\section{Réflexion critique sur l'Agentic AI}

Ce projet confirme les conclusions de notre analyse comparative : \textbf{l'Agentic
AI sans structuration SMA est fragile}. Les premières versions du prototype produisaient
des résultats incohérents. L'ajout progressif de patterns formels SMA a drastiquement
amélioré la fiabilité.

L'Agentic AI excelle pour les tâches cognitives (analyse sémantique, génération
naturelle). Elle échoue pour les tâches binaires précises (est-ce une offre
d'emploi ?). La solution : la compléter avec des algorithmes déterministes pour
les tâches où la précision prime sur la flexibilité --- c'est le paradigme
neuro-symbolique hybride identifié dans les études.

\section{Perspectives}

\paragraph{Kubernetes.}
Chaque service est sans état local (state dans ChromaDB et SQLite). Le scaling
horizontal d'A4 et A5 est directement applicable.

\paragraph{API LinkedIn officielle.}
Éliminerait la limitation principale et améliorerait significativement la précision
des scores.

\paragraph{MARL pour l'amélioration continue.}
Les agents A3a (Stratège) pourraient adapter leurs requêtes selon les scores
obtenus --- un apprentissage au sens strict.

\paragraph{Gouvernance et traçabilité.}
Un registre de décisions (quel agent, quelle confiance) adresserait les enjeux
de gouvernance identifiés par Milosevic et Rabhi (2026).

% =============================================================================
\chapter*{Références bibliographiques}
\addcontentsline{toc}{chapter}{Références bibliographiques}
% =============================================================================

\begin{description}[leftmargin=0pt, itemsep=4pt]

\item \textbf{Études SMA}

\item Dorri, A., Kanhere, S., Jurdak, R. (2018). \textit{Multi-Agent Systems: A Survey}. IEEE Access, 6, 28573--28593.

\item Maldonado, J. et al. (2024). \textit{Architectures and Coordination Patterns in Multi-Agent Systems}. Journal of AI Research, 79.

\item Nguyen, T. et al. (2019). \textit{Deep Reinforcement Learning for Multi-Agent Systems}. IEEE Transactions on Cybernetics, 50(9), 3826--3839.

\item General Survey (2020). \textit{A General Survey on Multi-Agent Coordination Mechanisms}. ACM Computing Surveys.

\item Gronauer, S., Diepold, K. (2021). \textit{Multi-Agent Deep Reinforcement Learning: A Survey}. AI Review, 55, 895--943.

\item Wang, J. et al. (2024). \textit{Multi-Agent Systems in Industry 4.0}. WJARR.

\medskip
\item \textbf{Études Agentic AI}

\item Wavestone (2024). \textit{Agentic AI Playbook --- Architecture Patterns for Enterprise Systems}.

\item Abou Ali, H. et al. (2026). \textit{Symbolic vs Neural Paradigms in Agentic AI Systems}. Expert Systems with Applications.

\item Alva, R., Pandey, S. (2024). \textit{Modular Architecture for LLM-Based Agentic Systems}. arXiv.

\item Sapkota, B. et al. (2024). \textit{AI Agents vs. Agentic AI: A Conceptual Taxonomy}. arXiv.

\item Patel, N. et al. (2025). \textit{Security Risks in LLM-Based Agentic Systems}. ACM CCS.

\item Hughes, A. et al. (2025). \textit{Safety and Alignment Challenges in Autonomous AI Agents}. Nature MI.

\item Vaidhyanathan, K., Taibi, D. (2026). \textit{Framework Comparison for Agentic AI}. Software: Practice \& Experience, 56(3).

\item TOSEM (2024). \textit{A Year in TOSEM --- SE4AI and AI4SE}.

\item Milosevic, D., Rabhi, F. (2026). \textit{Agentic Communities: Governance and Human-AI Collaboration}. IEEE Software.

\item Schneider, J. (2025). \textit{Agentic AI: Toward Goal-Directed Autonomous Systems}. IEEE Intelligent Systems.

\item Pati, S. et al. (2025). \textit{Agentic AI Paradigm: Planning, Memory and Adaptation}. AI Magazine.

\medskip
\item \textbf{Documents produits}

\item Dzermane, S. (2025). \textit{Synthèse Générale --- SMA et Éléments Connexes}. TER M1, Université de Montpellier.

\item Dzermane, S. (2025). \textit{Synthèse Générale --- Agentic AI et Éléments Connexes}. TER M1.

\item Dzermane, S. (2025). \textit{Analyse Comparative --- Agentic AI et Systèmes Multi-Agents}. TER M1.

\item Dzermane, S. (2025). \textit{Glossaire des concepts --- Comparaison des frameworks Agentic AI}. TER M1.

\end{description}

\end{document}
